% !TEX TS-program = xelatex
% !TEX encoding = UTF-8 Unicode
% !Mode:: "TeX:UTF-8"

\documentclass{resume}
\usepackage{zh_CN-Adobefonts_external} % Simplified Chinese Support using external fonts (./fonts/zh_CN-Adobe/)
%\usepackage{zh_CN-Adobefonts_internal} % Simplified Chinese Support using system fonts
\usepackage{linespacing_fix} % disable extra space before next section
\usepackage{soul}
\usepackage{cite}
\usepackage{hyperref} % for \href command
\hypersetup{
    colorlinks=false,   % 不使用颜色区分链接
    pdfborder={0 0 0},  % 去掉链接边框（三个参数分别对应左、右、下边框宽度）
}
\usepackage{fontawesome5} %add logo
\begin{document}
\pagenumbering{gobble} % suppress displaying page number

\name{聂小明}

\begin{center}
\faPhoneSquare \ (+86)138-0000-0000 \quad 
\href{mailto:example@domain.com}{\faEnvelopeO \ example@domain.com} \quad 
\href{https://github.com/USERNAME}{\faGithub @USERNAME}
\end{center}

% \section{\faGraduationCap\ 教育背景}
\section{\faGraduationCap \kaishu 教育背景}
% \begin{onehalfspacing}
\datedsubsection{\textbf{清华大学}，计算机科学与技术, \textit{在读硕士研究生}}{2025.9 - 2028.7}
\datedsubsection{\textbf{浙江大学}, 软件工程, \textit{工学学士}}{2021.9 - 2025.7}
\begin{itemize}
    \item \textbf{学业荣誉:} 专业排名1/100, GPA 3.9/4.0, \textbf{国家奖学金}, \textbf{优秀毕业生}, 一等奖奖学金
    \item \textbf{语言能力:} 英语 (CET-6), 日语 (JLPT N2)
\end{itemize}

% \end{onehalfspacing}

\section{\faBriefcase \kaishu 实习经历}
\datedsubsection{\textbf{字节跳动 | ByteDance} \kaishu 核心基础研究部}{2025.07--2025.12}
% \begin{onehalfspacing}
% === 项目 1：推荐系统性能优化 ===
\subsection{推荐排序系统在线性能提升}
\begin{itemize}
  \item \textbf{背景及问题}: 核心推荐排序链路在高并发场景下出现显著尾延迟，导致整体用户点击率下降。
  \item \textbf{工作与成果}: 设计并落地分布式缓存一致性方案，结合动态负载感知调度与线上灰度验证，将 99.9% 延迟缩短 40%，同时提升推荐召回准确率约 2%。
\end{itemize}
% === 项目 2：大规模日志平台架构升级 ===
\subsection{大规模日志平台架构升级}
\begin{itemize}
  \item \textbf{背景与问题}：原有日志采集系统在 PB 级数据下存在写入瓶颈，处理延迟高，导致线上监控数据滞后。
  \item \textbf{工作与成果}：参与设计基于 Flink 的流式数据处理框架，优化序列化与存储策略；实现端到端延迟降低 60%，支持 10+ 亿事件/天的稳定处理。
\end{itemize}
% \end{onehalfspacing}

\datedsubsection{\textbf{阿里巴巴 | Alibaba} \kaishu 机器智能与算法平台部}{2024.03--2024.07}
% \begin{onehalfspacing}
% === 项目 1：搜索排序模型加速 ===
\subsection{搜索排序模型在线加速与量化部署}
\begin{itemize}
  \item \textbf{背景及问题}: 搜索排序模型多层 Transformer 导致在线推理延迟高，影响搜索体验。
  \item \textbf{工作与成果}: 基于 ONNX Runtime 实现模型量化与分层缓存，开发端到端推理优化框架，使单请求延迟降低 35%，部署成本下降 25%。
\end{itemize}
% === 项目 2：智能推荐召回质量提升 ===
\subsection{智能推荐召回质量提升}
\begin{itemize}
  \item \textbf{背景及问题}：召回阶段多模态特征融合存在特征失衡，导致冷启动内容曝光不足。
  \item \textbf{工作与成果}：搭建多模态特征增强流水线，设计在线学习模块，实现召回覆盖率提升 18%，冷启动命中率提高 22%。
\end{itemize}
% \end{onehalfspacing}

\section{\faGithub \kaishu 开源经历}
% \begin{onehalfspacing}
\begin{itemize}
  \item 参与腾讯“犀牛鸟计划”开源项目孵化，主导多模态检索引擎模块设计与性能优化。
  \item 在“开放原子开源大赛”中提交并维护基于 Kubernetes 的边缘计算调度组件，入围决赛。
  \item 参与“开源之夏”社区活动，协作改进某主流深度学习框架的 C++ 算子性能，贡献 PR 若干。
\end{itemize}
% \end{onehalfspacing}

\section{\faMedal \kaishu 竞赛获奖}
% increase linespacing [parsep=0.5ex]
\datedsubsection{\textbf{全国大学生人工智能创新大赛,} 一等奖}{2026.02}
\datedsubsection{\textbf{某厂算法挑战赛,} 冠军}{2025.12}
\datedsubsection{\textbf{全国大学生创新大赛,} 一等奖}{2025.11}
\datedsubsection{\textbf{Kaggle竞赛,} 金牌}{2025.11}


\section{\faFilesO \kaishu 科研经历}
\begin{enumerate}[label={[\arabic*]}, leftmargin=*, itemsep=0pt, parsep=0pt]
  \item
  \textbf{Nie Xiaoming}, Xiaowen Liu, Jingfei Song, Rui Wang, and Hao Li*.
  \href{https://example.com}{\ul{Your Paper Name1}}.
  In \textit{International Conference on Systems and Machine Learning (SysML)}, 2025.

  \item
  Junjie Sun, Lin Mei*, \textbf{Nie Xiaoming}, Yifan Zhou, Qian Liu, Ming Xu
  \href{https://example.com}{\ul{Your Paper Name2}}.
  In \textit{Conference on Computer Vision and Pattern Recognition (CVPR)}, 2025.
\end{enumerate}

%% Reference
%\newpage
%\bibliographystyle{IEEETran}
%\bibliography{mycite}
\end{document}
